\chapter{Progressive implementation laboratories}

\section{Engineering principles}

The laboratory sequence moves from setup and repository preparation to isolated task execution, independent review and governance. Use disposable repositories with synthetic data. Run the offline laboratories first; those requiring provider credentials, MCP dependencies or GitHub administration are explicitly environment-dependent. No lab requires production access.

\section{Implementation approach}

The examples bundle contains the common application, tests, skills, task contract, orchestration code and policy. Reset a disposable copy for each lab, record command versions and retain a lab evidence folder outside the worker workspace. The laboratory table specifies a measurable pass condition; failures must be diagnosed rather than bypassed.

\section{Verification and limitations}

A lab passes only when the named evidence is collected and inspected. A documented but unexecuted command is not a verified result. Use the offline mock run to establish the mechanics before attempting live agents.

\subsection{Lab 01: Configure coding agents}

\textbf{Objective.} Configure coding agents using the reference architecture in Figure~\ref{fig:d23}.

\textbf{Prerequisites.} Installed approved CLIs; disposable checkout.

\textbf{Procedure.} Record versions; inspect help; run an interactive read-only repository summary in each.

\textbf{Evidence.} Version and permission evidence. \textbf{Pass criterion.} No write occurs; documented command accepted.

\textbf{Troubleshooting and safety.} Missing login: use approved account; never bypass.

\subsection{Lab 02: Agent-friendly repository}

\textbf{Objective.} Agent-friendly repository using the reference architecture in Figure~\ref{fig:d23}.

\textbf{Prerequisites.} Python 3.11+; example project.

\textbf{Procedure.} Inspect AGENTS.md, README, pricing.py and tests; run python -m unittest discover -s tests.

\textbf{Evidence.} Baseline test output. \textbf{Pass criterion.} All baseline assertions pass.

\textbf{Troubleshooting and safety.} If imports fail, run from project root.

\subsection{Lab 03: Repository instructions}

\textbf{Objective.} Repository instructions using the reference architecture in Figure~\ref{fig:d23}.

\textbf{Prerequisites.} Both CLIs for live check.

\textbf{Procedure.} Inspect AGENTS.md and CLAUDE.md import; request the approved test command.

\textbf{Evidence.} Instruction discovery record. \textbf{Pass criterion.} Both identify same command without executing unsafe actions.

\textbf{Troubleshooting and safety.} Confirm version and scope if instructions differ.

\subsection{Lab 04: Reusable Agent Skills}

\textbf{Objective.} Reusable Agent Skills using the reference architecture in Figure~\ref{fig:d23}.

\textbf{Prerequisites.} Python; skill packages.

\textbf{Procedure.} Run install\_skills.py; review metadata and trigger descriptions; compare copies.

\textbf{Evidence.} Seven installed skill packages per platform. \textbf{Pass criterion.} All files match canonical source.

\textbf{Troubleshooting and safety.} Remove stale copies before reinstalling.

\subsection{Lab 05: Connect MCP}

\textbf{Objective.} Connect MCP using the reference architecture in Figure~\ref{fig:d23}.

\textbf{Prerequisites.} MCP SDK and approved CLI.

\textbf{Procedure.} Install reviewed SDK in a venv; run policy\_mcp.py; register stdio server per platform docs; query allowed policy.

\textbf{Evidence.} Tool call and denied unknown ID. \textbf{Pass criterion.} Known ID returns policy; unknown ID rejected.

\textbf{Troubleshooting and safety.} Confirm transport and SDK version; no network secrets.

\subsection{Lab 06: Specification-driven feature}

\textbf{Objective.} Specification-driven feature using the reference architecture in Figure~\ref{fig:d23}.

\textbf{Prerequisites.} Toy pricing project.

\textbf{Procedure.} Read task.json; run trusted tests; execute mock or approved provider in disposable copy.

\textbf{Evidence.} Task contract and candidate diff. \textbf{Pass criterion.} All contract conditions and tests pass.

\textbf{Troubleshooting and safety.} Check base SHA and allowed paths.

\subsection{Lab 07: Verification loop}

\textbf{Objective.} Verification loop using the reference architecture in Figure~\ref{fig:d23}.

\textbf{Prerequisites.} Python; unit tests.

\textbf{Procedure.} Run orchestrate.py using mock; run unit tests for failed verifier and scope denial.

\textbf{Evidence.} JSON execution log. \textbf{Pass criterion.} Success and failure terminal states distinguishable.

\textbf{Troubleshooting and safety.} Do not interpret model text as success.

\subsection{Lab 08: Isolated Git worktrees}

\textbf{Objective.} Isolated Git worktrees using the reference architecture in Figure~\ref{fig:d23}.

\textbf{Prerequisites.} Git; committed disposable project.

\textbf{Procedure.} Run git worktree add -b task-a ../task-a HEAD; repeat for task-b; edit separately.

\textbf{Evidence.} Worktree list and separate diffs. \textbf{Pass criterion.} Changes remain in respective working trees.

\textbf{Troubleshooting and safety.} Use distinct paths; worktree is not sandbox.

\subsection{Lab 09: Agent review workflow}

\textbf{Objective.} Agent review workflow using the reference architecture in Figure~\ref{fig:d23}.

\textbf{Prerequisites.} Approved second CLI.

\textbf{Procedure.} Provide candidate diff to reviewer in separate read-only checkout; classify findings.

\textbf{Evidence.} Review report with file references. \textbf{Pass criterion.} No writes; every finding has evidence.

\textbf{Troubleshooting and safety.} Treat uncertain claims as hypotheses.

\subsection{Lab 10: GitHub Actions}

\textbf{Objective.} GitHub Actions using the reference architecture in Figure~\ref{fig:d23}.

\textbf{Prerequisites.} Disposable GitHub repo; admin.

\textbf{Procedure.} Review example workflow; configure trusted verifier and required status checks.

\textbf{Evidence.} CI run and rejected merge attempt. \textbf{Pass criterion.} Failing candidate cannot merge.

\textbf{Troubleshooting and safety.} Never enable secrets in untrusted PR job.

\subsection{Lab 11: Security remediation}

\textbf{Objective.} Security remediation using the reference architecture in Figure~\ref{fig:d23}.

\textbf{Prerequisites.} Toy project; approved scanner.

\textbf{Procedure.} Run negative tests; record finding; use remediation skill; rerun tests and scanner.

\textbf{Evidence.} Finding record with before/after evidence. \textbf{Pass criterion.} Original invalid input rejected; no suppression.

\textbf{Troubleshooting and safety.} Scanner unavailable: mark unverified.

\subsection{Lab 12: Hybrid workflow}

\textbf{Objective.} Hybrid workflow using the reference architecture in Figure~\ref{fig:d23}.

\textbf{Prerequisites.} Both CLIs; isolated verifier.

\textbf{Procedure.} Use one provider to implement; pass diff to other in read-only environment; run trusted gate.

\textbf{Evidence.} Implementation, review, test evidence. \textbf{Pass criterion.} Independent checks pass; review is advisory.

\textbf{Troubleshooting and safety.} Do not give reviewer push rights.

\subsection{Lab 13: Bounded multi-agent workflow}

\textbf{Objective.} Bounded multi-agent workflow using the reference architecture in Figure~\ref{fig:d23}.

\textbf{Prerequisites.} Git; two isolated tasks.

\textbf{Procedure.} Allocate disjoint scopes and worktrees; set retry budget; integrate after checks.

\textbf{Evidence.} Dependency graph and task reports. \textbf{Pass criterion.} No unauthorized overlapping edits.

\textbf{Troubleshooting and safety.} Escalate conflict rather than force merge.

\subsection{Lab 14: Evaluation and performance}

\textbf{Objective.} Evaluation and performance using the reference architecture in Figure~\ref{fig:d23}.

\textbf{Prerequisites.} Synthetic or measured JSONL records.

\textbf{Procedure.} Run evaluate.py records.jsonl; examine missing cost telemetry.

\textbf{Evidence.} Metric JSON. \textbf{Pass criterion.} Denominator includes failures and unknowns.

\textbf{Troubleshooting and safety.} Do not treat synthetic records as benchmark.

\subsection{Lab 15: Enterprise-ready environment}

\textbf{Objective.} Enterprise-ready environment using the reference architecture in Figure~\ref{fig:d23}.

\textbf{Prerequisites.} Sandbox and governance owners.

\textbf{Procedure.} Review threat matrix, ADR and runbook; test denied path, egress and release authorization.

\textbf{Evidence.} Control evidence and residual-risk register. \textbf{Pass criterion.} All mandatory controls validated by owners.

\textbf{Troubleshooting and safety.} Block deployment when evidence missing.

\section{Chapter conclusion}

Apply the stated verification criteria before expanding the scope of this workflow. Record deviations, evidence and unresolved risks in the task record.
